\section{Introduction}

Managed freight lanes are the primary capacity investment mechanism for Tier 1 Primary Arteries under Interstate 2.0. The distinction between managed freight lanes and general purpose lane additions is not merely operational but structural: general purpose lane additions induce demand and fill to capacity within 10 years \citep{Duranton2011}, while dedicated freight-only lanes with access control maintain their capacity benefit permanently by design. The fundamental problem with general purpose lane additions is that they add capacity available to all vehicle classes; induced demand fills the new capacity with the same mix of freight and passenger vehicles as the original lanes, restoring the original V/C ratio at higher total volume. Managed freight lanes, by restricting access to freight-classified vehicles, prevent induced passenger demand from consuming the new capacity.

This is not a marginal technical distinction. It is the difference between a 10-year improvement and a permanent infrastructure asset. The \$177--265 billion managed freight lane program for all 8 T1 corridors represents a capital allocation comparable to the original interstate highway program at 2024 dollars. The investment is justified only if the capacity benefit is permanent --- and permanence requires access control, not merely designation.

\subsection*{The Investment Case}

Two dedicated freight lanes per direction on I-75 Atlanta add 45,600 vehicle-per-day (vpd) of dedicated freight capacity to the corridor. The computation follows Highway Capacity Manual standards \citep{HCM7}: 2 lanes $\times$ 1,900 passenger car equivalents per hour per lane (pcphpl) $\times$ 24 hours = 91,200 vehicle-equivalent capacity additions per day, divided by a truck passenger car equivalency factor of 2.0 on level terrain, giving approximately 45,600 commercial vehicle equivalents per day of net new freight capacity.

This capacity addition produces four compounding benefits that together constitute the managed lane investment case:

\textbf{1. Dedicated capacity.} Freight lanes maintain V/C $\leq$ 0.85 (LOS C--D) regardless of induced demand growth in general purpose lanes, because induced demand accrues to the general purpose lanes, not the access-controlled freight lanes. The freight V/C is a function of freight demand growth only, which grows at 1.8\% per year per FHWA projections \citep{FHWA_freight2023}, compared to general purpose demand growth of 2.4\% per year.

\textbf{2. GP lane LOS improvement.} Removing freight vehicles from general purpose lanes improves passenger car LOS on existing GP lanes without adding GP lane capacity. On I-75 Atlanta, commercial vehicles represent approximately 22\% of total AADT. Transferring this 22\% to dedicated freight lanes reduces effective GP lane V/C from its current peak of 1.8+ to a projected 1.4, still congested but below the breakdown threshold that triggers cascading stop-and-go conditions.

\textbf{3. Sustained speed and platooning enablement.} Access-controlled freight lanes enable sustained 65 mph operation without passenger vehicle interference and enable truck platooning, which reduces aerodynamic drag by 20--35\% for following vehicles \citep{Browand2004}. At current diesel prices and mileage, 25\% fuel savings on I-75 Atlanta freight operations reduces shipper fuel cost by approximately \$0.08 per mile --- an annual savings of \$180 million across the corridor's freight volume.

\textbf{4. PTI improvement and shipper SLA capability.} The Planning Time Index (PTI) on T1 corridors currently ranges from 1.8 to 2.2, meaning shippers must plan for arrival windows 80--120\% longer than free-flow transit time to achieve 95th-percentile on-time delivery. Managed freight lanes target PTI $\leq$ 1.15 for the dedicated lanes, enabling shippers to offer tighter delivery windows --- an economic benefit captured in higher freight rates and reduced inventory buffer requirements.

\subsection*{Program Scope and Cost}

The full managed freight lane program covers the 7 T1 corridors with V/C above the I2.0 target (I-40 is excluded; see Section~\ref{sec:throughput}): I-10, I-75, I-80, I-90, I-95, I-85, and I-35. At \$10--15 million per lane-mile (the FHWA estimate for managed lane construction in constrained urban corridors \citep{FHWA_ML2021}), the program costs \$177--265 billion total across all T1 corridors. This is substantially less than the original interstate highway system cost at 2024 dollars (\$530 billion equivalent), reflecting the infrastructure advantage of building within existing right-of-way rather than greenfield.

\subsection*{Contributions}

\begin{enumerate}
\item Throughput analysis by corridor: V/C ratios, PTI measurements, and managed lane capacity addition for all 8 T1 corridors using HPMS data
\item Transit time modeling: NY--LA current and I2.0 managed-lane transit times, accounting for HOS regulations and freight platooning
\item NPV analysis: annualized construction cost vs. freight benefit (ATRI cost per truck-hour, V/C improvement, PTI reduction) on highest-volume segments
\item The I-40 exception: identification of the one T1 corridor where managed lanes are not needed, validating the tier classification methodology
\end{enumerate}
